\chapter{Objetivos del Proyecto}

\section{Objetivo general}
Desarrollar e implementar una herramienta gráfica interactiva para la codificación y decodificación de imágenes en escala de grises utilizando un autómata celular basado en la Regla 90 modificada, con el fin de explorar su comportamiento como técnica de transformación reversible de datos visuales.

\section{Objetivo del inciso uno}
Convertir imágenes en escala de grises a un arreglo binario que pueda ser procesado por un autómata celular, asegurando la integridad de la estructura original para su posterior reconstrucción.

\section{Objetivo inciso dos}
Aplicar la Regla 90 modificada de autómatas celulares a los datos binarios, permitiendo la transformación controlada de la imagen mediante un número configurable de iteraciones.

\section{Objetivo inciso tres}
Desarrollar un sistema de decodificación capaz de revertir con precisión el proceso de codificación, restaurando la imagen original a partir de los datos binarios transformados, siempre que se utilice el mismo número de iteraciones.

\section{Objetivo inciso cuatro}
Diseñar e implementar una interfaz gráfica de usuario que facilite la selección de archivos, configuración de iteraciones, visualización de resultados y manejo de errores de forma intuitiva y accesible.

\section{Objetivo inciso cinco}
Decodificar correctamente la imagen proporcionada por otros equipos, asegurando que el resultado coincida con la imagen original que dichos equipos codificaron previamente
